\section{Conclusion}
\label{sec:conclusion}

This paper has developed a two-region dynamic framework for evaluating public co-creation hubs in thin regional ecosystems. The model distinguishes between two margins of regional policy. The first is \emph{local coordination}: a public hub raises the rate at which potentially valuable local interactions are realized within a target region. The second is \emph{cross-regional renewal}: a complementary pipeline policy creates outside linkages that bring additional collaborators, ideas, and opportunities into the regional system. The paper's central claim is that these two margins should not be evaluated separately.

Three conclusions follow. First, a hub is socially efficient as a stand-alone intervention only if its local coordination gains are large enough to offset fixed cost and the dynamic redundancy created by repeated local interaction. The welfare case for a hub therefore depends not only on ecosystem thickness and the hub's matching effect, but also on the pace at which novelty is exhausted inside the local network. Second, short-run success is not equivalent to social efficiency. A hub may generate visible early activity and still be socially undesirable once repeated local ties and dynamic novelty loss are taken into account. Third, the relevant policy comparison is often not between a hub and no hub, but between a hub alone and a hub embedded in a broader regional strategy. A stand-alone hub may fail a welfare test even though a hub combined with cross-regional pipeline formation is socially efficient.

The extensions reinforce these conclusions. Participant composition matters in addition to headcount, so the value of a hub depends not only on how many actors a region contains, but also on whether those actors include highly absorptive or boundary-spanning participants. Conversely, digital coordination technologies weaken the case for physical hub creation when they substitute for baseline local matching, although they do not by themselves eliminate the problem of dynamic redundancy within a narrow regional network.

The paper also points to a more demanding empirical agenda. A theory-consistent evaluation of hub policy should not rely only on attendance, occupancy, or other utilization-based indicators. It should also examine whether new ties continue to form, whether repeated collaborator ratios rise over time, whether outside-region collaboration expands, and whether regions with stronger absorptive capacity derive larger gains from the same policy design. In this sense, the paper suggests that empirical work on hubs should be dynamic, network-aware, and explicitly regional.

The framework is intentionally stylized, and several extensions remain open. A richer model could allow for endogenous participant entry, repeated periods beyond the current two-period setting, strategic interaction across jurisdictions, or explicit network formation rather than reduced-form matching. An empirical extension could bring the theory closer to administrative or project-level data on collaboration networks, outside partners, and participant turnover.

Even in its simplified form, however, the model delivers a clear message for regional policy. In thin regional ecosystems, public hubs should not be judged as stand-alone venues that either succeed or fail according to visible early activity. They should be evaluated as components of a broader place-based strategy that combines local buzz with cross-regional renewal.
